\documentclass[10pt,twocolumn]{article}

% ─── Packages ───────────────────────────────────────────────────────────────
\usepackage[margin=0.75in, top=0.7in, bottom=0.7in]{geometry}
\usepackage{titlesec}
\usepackage{enumitem}
\usepackage{hyperref}
\usepackage{booktabs}
\usepackage{array}
\usepackage{microtype}
\usepackage{parskip}
\usepackage{xcolor}
\usepackage{fancyhdr}
\usepackage{lipsum}
\usepackage{amsmath}
\usepackage{amssymb}

% ─── Fix justified spacing in narrow columns ────────────────────────────────
\tolerance=1500
\emergencystretch=1em
\hyphenpenalty=50
\sloppy

% ─── Colors ──────────────────────────────────────────────────────────────────
\definecolor{accentblue}{RGB}{30, 80, 160}
\definecolor{lightgray}{RGB}{245, 245, 245}

% ─── Section formatting ──────────────────────────────────────────────────────
\titleformat{\section}
  {\normalfont\bfseries\color{accentblue}}
  {\thesection.}{0.5em}{}[\titlerule]

\titlespacing{\section}{0pt}{6pt}{3pt}

% ─── Hyperlinks ──────────────────────────────────────────────────────────────
\hypersetup{
  colorlinks=true,
  linkcolor=accentblue,
  urlcolor=accentblue,
  citecolor=accentblue
}

% ─── List spacing ────────────────────────────────────────────────────────────
\setlist[itemize]{leftmargin=*, nosep, topsep=2pt, itemsep=1pt}
\setlist[enumerate]{leftmargin=*, nosep, topsep=2pt, itemsep=1pt}

% ─── Header / Footer ─────────────────────────────────────────────────────────
\pagestyle{fancy}
\fancyhf{}
\renewcommand{\headrulewidth}{0.4pt}
\fancyhead[L]{\small\textcolor{accentblue}{\textbf{Project Proposal \& Progress Report}}}
\fancyhead[R]{\small Quantum Computing Course}
\fancyfoot[C]{\small\thepage}

% ─── Paragraph spacing ───────────────────────────────────────────────────────
\setlength{\parskip}{3pt}
\setlength{\parindent}{0pt}

% ════════════════════════════════════════════════════════════════════════════
\begin{document}

% ─── Title block (single column) ─────────────────────────────────────────────
\twocolumn[{%
  \begin{center}
    {\large\bfseries\color{accentblue}
      Replication and Barren Plateau Analysis of\\[2pt]
      Quantum Long Short-Term Memory Networks}\\[6pt]
    {\normalsize
      \textbf{Gautam Gupta} \quad
      \href{mailto:gautam23220@iiitd.ac.in}{\texttt{gautam23220@iiitd.ac.in}} \quad
      Roll No.: \textbf{2023220}}\\[4pt]
    {\small Indraprastha Institute of Information Technology Delhi (IIIT-Delhi)}
  \end{center}
  \vspace{4pt}
  \hrule
  \vspace{8pt}
}]


% ════════════════════════════════════════════════════════════════════════════
\section{Introduction}

Quantum Machine Learning (QML) has emerged as a compelling intersection of quantum computing and classical machine learning, leveraging superposition and entanglement for pattern modelling. Among the proposed hybrid architectures, the \textbf{Quantum Long Short-Term Memory (QLSTM)} network~\cite{chen2020qlstm} replaces parametric linear transformations inside classical LSTM cells with \emph{Variational Quantum Circuits} (VQCs), offering theoretical expressibility advantages for sequence modelling.

However, a fundamental and underexplored challenge in all VQC-based models is the \textbf{Barren Plateau (BP) phenomenon}. The seminal survey by Larocca et al.~\cite{larocca2024review} establishes that \emph{every component} of a variational quantum algorithm---the ansatz, initial state, observable, loss function, and hardware noise---can independently cause BPs, where the optimisation landscape becomes exponentially flat and featureless as system size grows. This is not merely a theoretical concern: on a BP, an exponentially large number of measurement shots is required to identify any loss-minimising direction, making training infeasible in practice. Despite this, the trainability of QLSTM under realistic circuit configurations has not been rigorously characterised.

This project closes that gap. I perform a faithful replication of the QLSTM architecture and then conduct a systematic empirical study of BP onset within its VQC sub-circuits. Crucially, I evaluate and compare three concrete mitigation strategies, including the recent two-stage convex-initialisation framework of Boabang \& Gyamerah~\cite{boabang2026twostage}, which shapes the quantum energy landscape into a smooth basin before nonconvex refinement.

The practical motivation is concrete: before investing resources in QLSTM for real-world tasks like \textbf{stock price forecasting} and \textbf{weather prediction}, a practitioner needs an evidence-based answer to---\emph{does this model actually train, and under what circuit configurations?}


% ════════════════════════════════════════════════════════════════════════════
\section{Project Description}

The project is structured in two tightly coupled phases.

\subsubsection*{Phase 1 — Faithful Replication}
I will implement the QLSTM architecture from Chen et al.~\cite{chen2020qlstm} using \texttt{PennyLane} and \texttt{PyTorch}. Each LSTM gate's linear transformation is replaced by a VQC sub-circuit using an angle-embedding layer followed by a strongly entangling ansatz. The primary benchmark is \textbf{stock price prediction} using daily closing prices from Yahoo Finance (e.g., S\&P~500 index). Weather forecasting (Delhi or Jena Climate dataset) is included as a stretch goal if time permits. A classical LSTM baseline with a \emph{matched parameter count} is trained for fair comparison on the primary benchmark.

\subsubsection*{Phase 2 — Barren Plateau Analysis \& Mitigation}
Guided by the taxonomy in Larocca et al.~\cite{larocca2024review}, I identify the specific BP sources active in QLSTM (ansatz expressibility and random initialisation) and measure \textbf{gradient variance} across VQC sub-circuits as a function of:
\begin{itemize}
  \item Circuit depth: $L \in \{1, 2, 4\}$
  \item Qubit count: $n \in \{2, 4, 6\}$
\end{itemize}
Gradient variance is estimated over 200 random parameter initialisations per configuration. I will then implement and compare \textbf{two} mitigation strategies:
\begin{enumerate}
  \item \textbf{Identity-block initialisation} (Grant et al.~\cite{grant2019init}) — sub-circuits initialised so the net unitary is close to identity, keeping initial gradients non-negligible.
  \item \textbf{Two-stage least squares} (Boabang \& Gyamerah~\cite{boabang2026twostage}) — a convex initialisation stage shapes the loss landscape into a smooth low-energy basin, followed by nonconvex refinement. This is the most recent (2026) mitigation framework in the literature.
\end{enumerate}
AI coding agents (as permitted by the course) will be used to accelerate boilerplate circuit construction and experiment scripting.

\subsubsection*{Papers to Implement / Survey}
\begin{enumerate}
  \item Chen et al.\ (2020). \emph{Quantum Long Short-Term Memory}. \href{https://arxiv.org/abs/2009.01783}{arXiv:2009.01783} --- primary architecture to replicate.
  \item Larocca et al.\ (2024). \emph{Barren Plateaus in Variational Quantum Computing: A Review}. \href{https://arxiv.org/abs/2405.00781}{arXiv:2405.00781} --- unified theoretical framework for BP sources and mitigations; primary reference for Phase 2 design.
  \item Boabang \& Gyamerah (2026). \emph{Overcoming Barren Plateaus via a Two-Step Least Squares Approach}. \href{https://arxiv.org/abs/2601.18060}{arXiv:2601.18060} --- mitigation strategy 2; novel two-stage optimisation framework.
\end{enumerate}


% ════════════════════════════════════════════════════════════════════════════
\section{Expected Outcome}

By the end of the project, I will produce:
\begin{itemize}
  \item A working \texttt{PennyLane}/\texttt{PyTorch} QLSTM implementation reproducing the key results of Chen et al.\ on the stock price benchmark.
  \item Gradient variance plots across circuit depths $L \in \{1,2,4\}$ and qubit counts $n \in \{2,4,6\}$, characterising BP onset and cross-referenced with the taxonomy of Larocca et al.~\cite{larocca2024review}.
  \item A comparison of \textbf{two} mitigation strategies (identity-block initialisation and two-stage least squares) with training loss curves and final MAE/RMSE.
  \item A quantitative QLSTM vs.\ classical LSTM comparison table: prediction error, parameter count, and training time.
  \item A clear, honest assessment of which circuit configurations make QLSTM trainable, grounded in BP theory.
\end{itemize}


% ════════════════════════════════════════════════════════════════════════════
\section{Midterm Deliverables}
% \textit{(Target: $\sim$15 hours of work; in-class presentation + report)}

\begin{itemize}
  \item Reading notes for all three papers, with a written summary of the BP taxonomy from Larocca et al.
  \item Working QLSTM implementation that trains on a toy dataset (sine wave) confirms circuit wiring and gradient flow are correct.
  \item Classical LSTM baseline with matched parameter count on the stock price dataset.
  \item Preliminary gradient variance measurements at $n \in \{2, 4\}$ qubits and $L \in \{1, 2\}$ layers, confirming the experimental setup is reproducible.
\end{itemize}


% ════════════════════════════════════════════════════════════════════════════
\section{Final Deliverables}
% \textit{(Target: $\sim$20 hours of work; viva/demo + presentation + report)}

\begin{itemize}
  \item Full QLSTM replication on the stock price dataset with a classical LSTM comparison table (MAE, RMSE, parameter count, training time). Weather forecasting included as a stretch goal.
  \item Complete BP analysis: gradient variance plots across $L \in \{1,2,4\}$ and $n \in \{2,4,6\}$, annotated using BP-source taxonomy from Larocca et al.
  \item Comparison of both mitigation strategies (identity-block initialisation and two-stage least squares) with training loss curves and final accuracy numbers.
  \item Documented code repository with reproducible scripts, ready for demo during viva.
  \item Final report and presentation summarising all results and conclusions.
\end{itemize}


% ════════════════════════════════════════════════════════════════════════════
% ════════════════════════════════════════════════════════════════════════════
%                        PROGRESS REPORT
% ════════════════════════════════════════════════════════════════════════════
% ════════════════════════════════════════════════════════════════════════════

\vspace{6pt}
\begin{center}
\rule{0.8\linewidth}{0.8pt}\\[4pt]
{\large\bfseries\color{accentblue} Progress Report}\\[2pt]
\rule{0.8\linewidth}{0.8pt}
\end{center}
\vspace{4pt}


% ════════════════════════════════════════════════════════════════════════════
\section{Papers Read}

\subsubsection*{1.\ Quantum Long Short-Term Memory}
\textbf{Authors:} Samuel Yen-Chi Chen, Shinjae Yoo, Yao-Lung L.\ Fang\\
\textbf{Venue:} arXiv preprint, arXiv:2009.01783 (2020)

The paper proposes QLSTM by replacing each linear transformation in the four LSTM gates (forget, input, cell, output) with a Variational Quantum Circuit (VQC). Each VQC first prepares qubits from the ground state using angle embedding (analogous to QFT/Hadamard for state preparation), then applies parameterised strongly entangling layers. Gradient computation uses the parameter shift rule ($\pm\pi/2$ shifts), which integrates with classical backpropagation via the chain rule to optimise the full hybrid model end-to-end.

\subsubsection*{2.\ Barren Plateaus in Variational Quantum Computing: A Review}
\textbf{Authors:} Martin Larocca, Supanut Thanasilp, Samson Wang, Kunal Sharma, Jacob Biamonte, Patrick J.\ Coles, Lukasz Cincio, Jarrod R.\ McClean, Zo\"{e} Holmes, M.\ Cerezo\\
\textbf{Venue:} arXiv preprint, arXiv:2405.00781 (2024)

This comprehensive review establishes a unified taxonomy of barren plateau (BP) sources in variational quantum algorithms. The core result is that the variance of the loss gradient $\text{Var}[\partial L/\partial\theta]$ shrinks exponentially as $O(1/2^n)$ with qubit count $n$, making gradient-based optimisation infeasible for large circuits. BPs arise independently from ansatz expressibility, entanglement structure, noise, and choice of global vs.\ local cost functions.

\subsubsection*{3.\ Overcoming Barren Plateaus via a Two-Step Least Squares Approach}
\textbf{Authors:} Francis Boabang, Samuel A.\ Gyamerah\\
\textbf{Venue:} arXiv preprint, arXiv:2601.18060 (2026)

Proposes a two-stage optimisation framework to escape barren plateaus. Stage~1 uses regularised least squares (a convex problem) to find a good initial parameter region, shaping the loss landscape into a smooth basin. Stage~2 refines these parameters using unregularised gradient-based optimisation. This avoids the problem of random initialisation landing on exponentially flat regions of the landscape.

\subsubsection*{4.\ Barren Plateaus in Quantum Neural Network Training Landscapes}
\textbf{Authors:} Jarrod R.\ McClean, Sergio Boixo, Vadim N.\ Smelyanskiy, Ryan Babbush, Hartmut Neven\\
\textbf{Venue:} \emph{Nature Communications}, 9(1):4812 (2018)

Seminal paper providing the first rigorous proof that for sufficiently random parameterised quantum circuits, gradient magnitudes vanish exponentially with qubit count due to concentration of measure. Establishes the mathematical foundation showing why deep, randomly initialised VQCs become untrainable, the result that motivates all subsequent BP mitigation research.

\subsubsection*{5.\ An Initialization Strategy for Addressing Barren Plateaus}
\textbf{Authors:} Edward Grant, Leonard Wossnig, Mateusz Ostaszewski, Marcello Benedetti\\
\textbf{Venue:} \emph{Quantum}, 3:214 (2019)

Proposes initialising VQC parameters such that each circuit block approximates the identity operation, analogous to classical He/Xavier initialisation. This ensures the circuit starts in a non-flat region of the optimisation landscape with non-vanishing gradients, enabling successful training before the circuit explores more expressive regions.


% ════════════════════════════════════════════════════════════════════════════
\section{Summary of Key Ideas}

The papers collectively address the challenge of training variational quantum circuits for machine learning. Three key themes emerge:

\begin{enumerate}
  \item \textbf{Hybrid Quantum-Classical Architectures:} Chen et al.\ demonstrate that VQCs can replace linear layers inside classical neural networks (specifically LSTMs). The parameter shift rule enables gradient computation through quantum circuits, which integrates with classical backpropagation via the chain rule---the $\pm\pi/2$ shift computes $\partial\langle O\rangle/\partial\theta$ (the quantum part), and PyTorch autograd handles $\partial L/\partial\langle O\rangle$ (the classical loss part).

  \item \textbf{The Barren Plateau Problem:} McClean et al.\ (2018) and Larocca et al.\ (2024) establish that randomly initialised VQCs suffer from exponentially vanishing gradients as circuit size grows. This is a fundamental obstacle: simply scaling up quantum circuits does not improve performance without careful design. BP sources include circuit expressibility (too expressive $\Rightarrow$ approximate random unitary $\Rightarrow$ flat landscape), entanglement, hardware noise, and global cost functions.

  \item \textbf{Mitigation Through Initialisation:} Both Grant et al.\ and Boabang \& Gyamerah address BPs through smarter initialisation. Grant et al.\ use identity-block initialisation (keeping initial unitaries near identity, analogous to classical He/Xavier), while Boabang \& Gyamerah use a two-stage approach where convex optimisation first finds a good starting region. Both ensure the optimiser begins where gradients are meaningful.
\end{enumerate}


% ════════════════════════════════════════════════════════════════════════════
\section{Implementation Status}

The following components have been implemented using \texttt{PennyLane}~(v0.42) and \texttt{PyTorch}~(v2.6):

\begin{enumerate}
  \item \textbf{QLSTM Architecture} (\texttt{src/qlstm.py}) — VQC sub-circuits with angle embedding + strongly entangling layers, QLSTM cell with four VQC-based gates, full model wrapper.
  \item \textbf{Classical LSTM Baseline} (\texttt{src/classical\_lstm.py}) — standard PyTorch LSTM for fair comparison.
  \item \textbf{Sine Wave Training} (\texttt{src/train\_sine.py}) — toy benchmark that confirms correct circuit wiring and gradient flow.
  \item \textbf{S\&P 500 Stock Training} (\texttt{src/train\_stock.py}) — daily closing-price prediction using \texttt{yfinance} data, training all three QLSTM variants and the classical baseline.
  \item \textbf{Barren Plateau Analysis} (\texttt{src/bp\_full\_grid.py}) — gradient variance over 200 random initialisations for the full proposal grid $n \in \{2, 4, 6\}$, $L \in \{1, 2, 4\}$.
  \item \textbf{Mitigation Strategies} (\texttt{src/mitigations.py}) — identity-block initialisation (Grant et al.) and two-stage LS warm-up (Boabang \& Gyamerah).
  \item \textbf{Trainability Table} (\texttt{src/trainability\_table.py}) — automatic classifier mapping each $(n, L)$ to a trainability verdict.
\end{enumerate}

\smallskip
\textbf{Repository:} \url{https://github.com/gautam2905/IQC-Project}

% ------------------------------------------------------------------
\subsection*{Barren Plateau Characterisation}

Using 200 random initialisations per configuration and taking the average variance across all circuit parameters, we obtained the following for the VQC sub-circuit used inside QLSTM gates:

\begin{center}
\begin{tabular}{@{}cccc@{}}
\toprule
\textbf{Layers} $\backslash$ \textbf{Qubits} & $n{=}2$ & $n{=}4$ & $n{=}6$ \\
\midrule
$L{=}1$ & $8.82\!\times\!10^{-2}$ & $2.61\!\times\!10^{-2}$ & $5.58\!\times\!10^{-3}$ \\
$L{=}2$ & $6.57\!\times\!10^{-2}$ & $1.40\!\times\!10^{-2}$ & $8.53\!\times\!10^{-3}$ \\
$L{=}4$ & $8.00\!\times\!10^{-2}$ & $2.35\!\times\!10^{-2}$ & $4.70\!\times\!10^{-3}$ \\
\bottomrule
\end{tabular}
\end{center}

Going from $n{=}2$ to $n{=}6$ qubits, the gradient variance shrinks by ${\sim}15\times$ — consistent with the exponential $O(1/2^n)$ decay predicted by McClean et al.~\cite{mcclean2018barren} and Larocca et al.~\cite{larocca2024review}.

% ------------------------------------------------------------------
\subsection*{Trainability Verdict per Configuration}

Using thresholds (Trainable $\ge 10^{-2}$, Marginal $10^{-4}$--$10^{-2}$, Barren $< 10^{-4}$):

\begin{center}
\begin{tabular}{@{}cccc@{}}
\toprule
\textbf{Layers} $\backslash$ \textbf{Qubits} & $n{=}2$ & $n{=}4$ & $n{=}6$ \\
\midrule
$L{=}1$ & Trainable & Trainable & Marginal \\
$L{=}2$ & Trainable & Trainable & Marginal \\
$L{=}4$ & Trainable & Trainable & Marginal \\
\bottomrule
\end{tabular}
\end{center}

At $n{=}6$, gradient magnitudes are already in the ``marginal'' regime, foreshadowing a barren plateau for slightly larger circuits. Practitioners should target $n \le 4$ for reliable training without mitigation.

% ------------------------------------------------------------------
\subsection*{S\&P 500 Stock Price Results}

Training configuration: sequence length $10$, $80$ training samples, $20$ epochs, \texttt{Adam} (lr = $0.02$). QLSTM: $n_q{=}2$, $L{=}1$, hidden = $4$.

\begin{center}
\begin{tabular}{@{}lcccc@{}}
\toprule
\textbf{Model} & \textbf{MAE} & \textbf{RMSE} & \textbf{Params} & \textbf{Time (s)} \\
\midrule
Classical LSTM (h=8)             & 2.363 & 2.367 & 361 & 28.3 \\
Classical LSTM matched (h=4)     & 2.000 & 2.001 & 117 & 26.3 \\
QLSTM (random init)              & \textbf{1.242} & \textbf{1.249} & 125 & 528.0 \\
QLSTM (identity-block init)      & 1.246 & 1.253 & 125 & 523.1 \\
QLSTM (two-stage LS)             & 2.374 & 2.389 & 125 & 533.4 \\
\bottomrule
\end{tabular}
\end{center}

The parameter-matched comparison uses Classical LSTM with \texttt{hidden\_size=4} (117 params) against QLSTM with $n_q{=}2$, $L{=}1$ (125 params) — a near-identical parameter budget. Three observations:

\begin{itemize}
  \item \textbf{QLSTM beats the size-matched classical baseline by ${\sim}37\%$ in MAE} (1.242 vs 2.000) with essentially the same parameter count (125 vs 117). This is a strong result: at small circuit sizes, QLSTM's quantum expressibility provides a genuine advantage over an equivalently sized classical model on this non-stationary time series.
  \item \textbf{Identity-block initialisation behaves nearly identically to random init} at $n{=}2$. This is expected: barren plateaus are not yet severe at two qubits (Var $\sim 10^{-1}$), so smart initialisation provides little benefit in this regime. Its value should become apparent at larger $n$ where BP is a real obstacle.
  \item \textbf{Two-stage LS warm-up degrades generalisation.} Stage~1 finds an analytically optimal linear projection onto the current (random) quantum features; however, subsequent gradient descent perturbs those quantum features, and the combination overfits the training window of the non-stationary S\&P~500 series. This is an \emph{honest null result}: the convex warm-up may help \emph{optimisation speed} but can hurt \emph{generalisation} under distribution shift. A regularised variant of the warm-up is proposed as the novel extension.
\end{itemize}

% ------------------------------------------------------------------
\subsection*{Sine Wave Sanity Check}

The earlier toy benchmark (sine wave, 50 train samples, 30 epochs) confirmed correct QLSTM gradient flow: loss decreased monotonically from $0.52$ to $0.48$; classical LSTM reached $0.24$ on the same task with $\sim\!4\times$ more parameters.


% ════════════════════════════════════════════════════════════════════════════
\section{Refined Final Deliverable}

All core deliverables from the original proposal are now implemented; the remaining work is empirical scaling and a novel extension.

\subsection*{Completed}
\begin{enumerate}
  \item QLSTM replication with angle embedding + strongly entangling ansatz.
  \item Classical LSTM baseline with matched architecture.
  \item BP gradient-variance sweep across the proposal grid $n \in \{2,4,6\}$, $L \in \{1,2,4\}$.
  \item Identity-block initialisation and two-stage LS warm-up mitigations.
  \item Full S\&P~500 benchmark comparing Classical LSTM against three QLSTM variants, with MAE / RMSE / parameter count / training-time table.
  \item Trainability verdict table mapping each $(n, L)$ pair to Trainable / Marginal / Barren.
\end{enumerate}

\subsection*{Planned (remaining)}
\begin{enumerate}
  \item \textbf{Scaling the BP sweep} from $n{=}6$ to $n \approx 20$ qubits, producing a clean log-linear $O(1/2^n)$ curve. This is purely a compute-time investment on the classical simulator.
  \item \textbf{Extended stock training}: larger training windows, more epochs, and error bars across multiple seeds for the Classical~vs.~QLSTM comparison. The current small-window run (80 samples, 20 epochs) was chosen to keep quantum-simulation cost manageable during the midterm phase.
  \item \textbf{Weather forecasting stretch goal} (Delhi / Jena Climate) as a secondary benchmark if time permits.
  \item \textbf{Novel extension (hypothesis, scope-dependent):} our S\&P~500 experiments suggest that two-stage LS warm-up can \emph{hurt} generalisation under distribution shift — it fits the training window tightly but becomes brittle downstream. We plan to investigate a \emph{regularised variant} of the warm-up that keeps the good training-time properties (smooth basin for Stage~2) while explicitly preserving generalisation, e.g.\ by combining the LS warm-up with an early-stopping criterion or a regularisation term tied to the Fisher information of the classical readout. A positive result would constitute a small but concrete contribution beyond pure replication.
  \item \textbf{Viva-ready notebook} bundling the full pipeline for reproducible demo.
\end{enumerate}


% ════════════════════════════════════════════════════════════════════════════
% ════════════════════════════════════════════════════════════════════════════
%                          FINAL REPORT
% ════════════════════════════════════════════════════════════════════════════
% ════════════════════════════════════════════════════════════════════════════

\vspace{6pt}
\begin{center}
\rule{0.8\linewidth}{0.8pt}\\[4pt]
{\large\bfseries\color{accentblue} Final Report}\\[2pt]
{\small Date of Submission: \today}\\[2pt]
\rule{0.8\linewidth}{0.8pt}
\end{center}
\vspace{4pt}


% ════════════════════════════════════════════════════════════════════════════
\section{Problem Statement}

Variational Quantum Circuits (VQCs) have been proposed as drop-in replacements for the linear layers inside classical recurrent architectures. The Quantum Long Short-Term Memory (QLSTM) network of Chen et al.\ \cite{chen2020qlstm} is the most prominent such hybrid: it swaps all four LSTM gate transformations for VQCs and claims expressibility advantages on sequence tasks. However, VQC-based models are known to suffer from the \textbf{Barren Plateau (BP)} phenomenon, where the gradient of the loss function vanishes exponentially with the number of qubits, making training infeasible at scale.

This project addresses two tightly linked questions:
\begin{enumerate}
  \item \textbf{Does QLSTM actually train?} We faithfully replicate the architecture and benchmark it against a parameter-matched classical LSTM on S\&P~500 stock price prediction.
  \item \textbf{Under what circuit configurations does BP onset occur?} We characterise gradient variance across circuit depths $L \in \{1,2,4\}$ and qubit counts $n \in \{2,4,6\}$, and evaluate two mitigation strategies — identity-block initialisation (Grant et al.\ \cite{grant2019init}) and two-stage least-squares warm-up (Boabang \& Gyamerah \cite{boabang2026twostage}) — against the random-init baseline.
\end{enumerate}

The practical outcome is a concrete answer to a question any practitioner would ask before deploying QLSTM: \textit{which circuit configurations are trainable, and do mitigation strategies help?}


% ════════════════════════════════════════════════════════════════════════════
\section{Papers}

\begin{enumerate}

\item \textbf{Quantum Long Short-Term Memory}\\
S.\ Y.-C.\ Chen, S.\ Yoo, Y.-L.\ L.\ Fang\\
arXiv preprint, arXiv:2009.01783, 2020

\item \textbf{Barren Plateaus in Variational Quantum Computing: A Review}\\
M.\ Larocca, S.\ Thanasilp, S.\ Wang, K.\ Sharma, J.\ Biamonte, P.\ J.\ Coles,
L.\ Cincio, J.\ R.\ McClean, Z.\ Holmes, M.\ Cerezo\\
arXiv preprint, arXiv:2405.00781, 2024

\item \textbf{Overcoming Barren Plateaus in Variational Quantum Circuits Using a Two-Step Least Squares Approach}\\
F.\ Boabang, S.\ A.\ Gyamerah\\
arXiv preprint, arXiv:2601.18060, 2026

\item \textbf{Barren Plateaus in Quantum Neural Network Training Landscapes}\\
J.\ R.\ McClean, S.\ Boixo, V.\ N.\ Smelyanskiy, R.\ Babbush, H.\ Neven\\
\textit{Nature Communications}, 9(1):4812, 2018

\item \textbf{An Initialization Strategy for Addressing Barren Plateaus in Parametrized Quantum Circuits}\\
E.\ Grant, L.\ Wossnig, M.\ Ostaszewski, M.\ Benedetti\\
\textit{Quantum}, 3:214, 2019

\end{enumerate}


% ════════════════════════════════════════════════════════════════════════════
\section{Methodology}

\subsection*{Module 1 — QLSTM (\texttt{src/qlstm.py})}

Implements the architecture of Chen et al.\ using PennyLane and PyTorch. The key class hierarchy is:

\begin{itemize}
  \item \textbf{\texttt{VQC}}: A single variational quantum circuit. Takes a classical input vector, projects it to $n$ qubits via a trainable \texttt{Linear} pre-net, applies \texttt{AngleEmbedding} (Y-rotation encoding), then \texttt{StronglyEntanglingLayers} ($L$ layers of parameterised rotations + CNOT ring), and returns $n$ expectation values $\langle Z_i \rangle$ projected back to the required output dimension via a trainable \texttt{Linear} post-net.
  \item \textbf{\texttt{QLSTMCell}}: Instantiates four independent \texttt{VQC}s — one per LSTM gate (forget, input, cell candidate, output) — and reproduces the standard LSTM update equations with those VQCs in place of the classical weight matrices.
  \item \textbf{\texttt{QLSTM}}: Wraps the cell into a sequence model; iterates over the time dimension and produces a final prediction from the last hidden state.
\end{itemize}

Gradient computation uses PennyLane's \texttt{backprop} method (which applies the parameter shift rule internally) and integrates with PyTorch autograd via the chain rule, enabling end-to-end training with standard optimisers.

\subsection*{Module 2 — Classical LSTM (\texttt{src/classical\_lstm.py})}

Standard \texttt{torch.nn.LSTM} wrapper. Two instances are used in experiments: a larger model (\texttt{hidden\_size=8}, 361 params) and a \emph{parameter-matched} model (\texttt{hidden\_size=4}, 117 params) for a fair comparison against QLSTM (125 params).

\subsection*{Module 3 — Barren Plateau Analysis (\texttt{src/bp\_full\_grid.py})}

For each $(n, L)$ configuration in the proposal grid, runs 200 forward-backward passes with random parameter and input initialisations. Collects the full gradient tensor per pass and computes the per-parameter variance, then averages across all parameters. This \emph{mean-over-all-parameters} metric is a specific implementation choice: using only the first parameter gives structural zeros for single-layer circuits (a Z-rotation applied before a Z-measurement contributes no gradient), so the mean is more representative and robust.

\subsection*{Module 4 — Mitigation Strategies (\texttt{src/mitigations.py})}

Two functions:

\begin{itemize}
  \item \textbf{\texttt{identity\_block\_init}}: Generates a \texttt{(n\_layers, n\_qubits, 3)} weight tensor where consecutive layer pairs cancel — layer $2k{+}1$ is set to the negative of layer $2k$ (exact inverse rotation), plus a small Gaussian perturbation $\epsilon{=}0.05$ to break the exact saddle. Applied to every VQC in the QLSTM via \texttt{apply\_identity\_init\_to\_qlstm}.
  \item \textbf{\texttt{two\_stage\_ls\_warmup}}: Registers a forward hook on \texttt{model.fc\_out} to capture pre-projection features $\Phi \in \mathbb{R}^{N \times d}$ for the full training set in one frozen forward pass. Solves $w^* = (\Phi^\top\Phi + \lambda I)^{-1}\Phi^\top y$ analytically and copies $w^*$ into the final linear layer. Stage~2 is then ordinary gradient descent from this warm start.
\end{itemize}

\subsection*{Module 5 — S\&P 500 Data Loader (\texttt{src/data\_loader.py})}

Downloads daily closing prices via \texttt{yfinance} (2020--2024, ${\sim}1000$ trading days), normalises by mean and standard deviation, and converts to sliding-window sequences of length 10. Caches the raw price array to avoid repeated API calls.

\subsection*{Module 6 — Training \& Evaluation (\texttt{src/train\_stock.py}, \texttt{src/train\_sine.py})}

Unified training loop (Adam, configurable lr and epochs) with timed wall-clock measurement. Evaluates MAE, RMSE, and MSE on a held-out test split. Generates loss-curve and prediction-overlay plots automatically.

\subsection*{Module 7 — Trainability Table (\texttt{src/trainability\_table.py})}

Reads the BP grid JSON and classifies each $(n, L)$ cell using fixed variance thresholds: Trainable ($\ge 10^{-2}$), Marginal ($10^{-4}$--$10^{-2}$), Barren ($< 10^{-4}$). Exports both a JSON and a \LaTeX{} tabular snippet.


% ════════════════════════════════════════════════════════════════════════════
\section{Observations}

\subsection*{Barren Plateau Characterisation}

\begin{center}
\begin{tabular}{@{}cccc@{}}
\toprule
\textbf{Layers} $\backslash$ \textbf{Qubits} & $n{=}2$ & $n{=}4$ & $n{=}6$ \\
\midrule
$L{=}1$ & $8.82\!\times\!10^{-2}$ \textit{(T)} & $2.61\!\times\!10^{-2}$ \textit{(T)} & $5.58\!\times\!10^{-3}$ \textit{(M)} \\
$L{=}2$ & $6.57\!\times\!10^{-2}$ \textit{(T)} & $1.40\!\times\!10^{-2}$ \textit{(T)} & $8.53\!\times\!10^{-3}$ \textit{(M)} \\
$L{=}4$ & $8.00\!\times\!10^{-2}$ \textit{(T)} & $2.35\!\times\!10^{-2}$ \textit{(T)} & $4.70\!\times\!10^{-3}$ \textit{(M)} \\
\bottomrule
\end{tabular}\\[3pt]
\small T = Trainable ($\ge 10^{-2}$), M = Marginal ($10^{-4}$--$10^{-2}$)
\end{center}

Gradient variance drops ${\sim}15\times$ from $n{=}2$ to $n{=}6$, consistent with the $O(1/2^n)$ theoretical prediction. All configurations in the grid remain at least marginal; barren-plateau onset is expected around $n \approx 8$--$10$.

\subsection*{S\&P 500 Stock Price Benchmark}

\begin{center}
\begin{tabular}{@{}lcccc@{}}
\toprule
\textbf{Model} & \textbf{MAE} & \textbf{RMSE} & \textbf{Params} & \textbf{Time (s)} \\
\midrule
Classical LSTM (h=8)          & 2.363 & 2.367 & 361 & 28.3 \\
Classical LSTM matched (h=4)  & 2.000 & 2.001 & 117 & 26.3 \\
QLSTM (random init)           & \textbf{1.242} & \textbf{1.249} & 125 & 528.0 \\
QLSTM (identity-block init)   & 1.246 & 1.253 & 125 & 523.1 \\
QLSTM (two-stage LS)          & 2.374 & 2.389 & 125 & 533.4 \\
\bottomrule
\end{tabular}
\end{center}

Key takeaways:

\begin{itemize}
  \item QLSTM (random init, $n{=}2$, $L{=}1$) outperforms the parameter-matched classical LSTM by \textbf{37\% in MAE} (1.242 vs 2.000) with a nearly identical parameter budget (125 vs 117).
  \item Identity-block initialisation performs on par with random init at $n{=}2$. At this qubit count, BP is not yet active (Var $\sim 10^{-1}$), so the initialisation strategy makes no practical difference — which is exactly what theory predicts.
  \item The two-stage LS warm-up \emph{degrades} test-set performance (MAE 2.374 $\approx$ classical baseline). The Stage~1 solution tightly fits the training window's quantum features; when Stage~2 then adjusts the quantum parameters, the classical readout becomes misaligned with the new features, and the combination generalises poorly on this non-stationary series. This is an honest null result and points to a concrete open problem (see Future Direction).
\end{itemize}

\subsection*{Repository}

All code, plots, and this report are available at:\\
\url{https://github.com/gautam2905/IQC-Project}


% ════════════════════════════════════════════════════════════════════════════
\section{Future Direction}

\textbf{Regularised two-stage LS (proposed novel extension).} The null result on the LS warm-up reveals a fundamental tension: Stage~1 optimally fits the current quantum features, but Stage~2 moves those features, destroying the fit. One principled fix is to \emph{jointly} regularise the quantum parameters during Stage~2 to stay close to the feature distribution seen in Stage~1 — for example, by adding a Fisher-information penalty that penalises large changes in the circuit's output distribution. Alternatively, one could \emph{freeze the classical readout} after Stage~1 and only fine-tune the quantum parameters, keeping the linear projection anchored to the well-initialised region. Either variant would retain the smooth-basin benefit of the convex warm-up while avoiding the distribution-shift problem.

\textbf{Scaling the BP sweep to $n \approx 20$.} The current grid stops at $n{=}6$ (marginal regime). Extending to $n{=}10$--$20$ on the classical simulator would produce a clean log-linear $O(1/2^n)$ curve, making the exponential decay visually unambiguous and quantitatively verifiable against the theoretical slope. This is purely a compute investment but would produce a compelling empirical plot.

\textbf{Local vs global cost functions as a structural BP mitigation.} Larocca et al.\ show that global cost functions (measuring all qubits together) cause worse BPs than local ones (measuring one qubit at a time). QLSTM currently uses a global PauliZ sum. Replacing this with per-qubit measurements and aggregating classically is a low-effort architectural change that could substantially improve trainability at $n \ge 6$ without changing the circuit structure at all.

\textbf{Hardware noise analysis.} All experiments here use a noiseless classical simulator (\texttt{default.qubit}). On real hardware, depolarising noise is another independent source of barren plateaus (Larocca et al., Source~4). Repeating the BP sweep with PennyLane's \texttt{default.mixed} device and a configurable error rate would characterise how quickly noise destroys trainability, which is the practically relevant question for near-term quantum devices.

\textbf{Weather forecasting benchmark.} The Delhi and Jena Climate datasets were included as a stretch goal in the proposal. Running the same QLSTM vs classical LSTM comparison on a multi-variate meteorological time series would test whether the observed quantum advantage on S\&P~500 generalises beyond financial data, or is specific to the structure of that particular series.


% ─── References ──────────────────────────────────────────────────────────────
\begin{thebibliography}{9}
\small

\bibitem{chen2020qlstm}
S.\ Chen, S.\ Yoo, and Y.-L.\ Fang.
\newblock Quantum long short-term memory.
\newblock \emph{arXiv preprint arXiv:2009.01783}, 2020.
\newblock \url{https://arxiv.org/abs/2009.01783}

\bibitem{larocca2024review}
M.\ Larocca, S.\ Thanasilp, S.\ Wang, K.\ Sharma, J.\ Biamonte, P.\ J.\ Coles,
L.\ Cincio, J.\ R.\ McClean, Z.\ Holmes, and M.\ Cerezo.
\newblock Barren plateaus in variational quantum computing: A review.
\newblock \emph{arXiv preprint arXiv:2405.00781}, 2024.
\newblock \url{https://arxiv.org/abs/2405.00781}

\bibitem{mcclean2018barren}
J.\ R.\ McClean, S.\ Boixo, V.\ N.\ Smelyanskiy, R.\ Babbush, and H.\ Neven.
\newblock Barren plateaus in quantum neural network training landscapes.
\newblock \emph{Nature Communications}, 9(1):4812, 2018.

\bibitem{grant2019init}
E.\ Grant, L.\ Wossnig, M.\ Ostaszewski, and M.\ Benedetti.
\newblock An initialization strategy for addressing barren plateaus in parametrized quantum circuits.
\newblock \emph{Quantum}, 3:214, 2019.

\bibitem{boabang2026twostage}
F.\ Boabang and S.\ A.\ Gyamerah.
\newblock Overcoming barren plateaus in variational quantum circuits using a two-step least squares approach.
\newblock \emph{arXiv preprint arXiv:2601.18060}, 2026.
\newblock \url{https://arxiv.org/abs/2601.18060}

\end{thebibliography}

\end{document}